% ======================================================================
\section{Introduction}\label{sec:intro}
% ======================================================================
Surface missions to the Moon need to know where they are without GNSS. Wheel and visual
odometry~\cite{maimone2007mer} drift without bound, so a rover periodically needs an absolute fix
or a loop closure. One way to get it is to recognise that the current view shows a place that is
already in a reference set, either imagery from an earlier traverse or views rendered from orbital
digital elevation models (DEMs)~\cite{johnson2008trn,ceresoli2025optical}. This is
\emph{visual place recognition} (VPR), usually posed as image retrieval against a geo-tagged
database~\cite{lowry2015visual,schubert2023tutorial}.

Modern VPR systems reach high recall on street-level benchmarks by aggregating deep local
features into a compact global descriptor~\cite{arandjelovic2016netvlad,radenovic2019gem,berton2022cosplace,berton2023eigenplaces,izquierdo2024optimal},
optionally followed by local-feature re-ranking~\cite{hausler2021patch}. The lunar surface
removes almost every cue these systems rely on:
\begin{itemize}
\item \textbf{Illumination.} There is no atmosphere to scatter light, so shadows are black and
sharp. Over a lunar day the Sun's elevation and azimuth sweep through large ranges, and the
same place can look entirely different at two times (\cref{fig:conds}).
\item \textbf{Self-similarity.} Craters and boulders follow power-law size--frequency
distributions~\cite{neukum2001cratering,melosh1989impact}, so terrain is statistically repetitive
at every scale and individual features are weakly discriminative~\cite{torii2013visual}.
\item \textbf{Photometry.} Regolith is strongly non-Lambertian, with phase-dependent limb
darkening and an opposition surge~\cite{mcewen1991photometric,hapke2012theory}, which violates the
brightness-constancy assumptions built into terrestrial descriptors.
\item \textbf{Data.} There is no large, pose-annotated lunar VPR dataset to train on.
\end{itemize}

Two design questions follow. \emph{How should tokens be aggregated} when a large and variable part of
every image (black sky, featureless regolith, moving shadows) carries no place information? And
\emph{how should local evidence be used} when local features respond mostly to shadow edges that move
with the Sun? This paper answers both within one framework, \textbf{Galatea}, which casts each stage
of two-stage VPR as an optimal-transport (OT) problem. Global aggregation becomes an assignment of
tokens to clusters and a dust-bin; local verification becomes a transport between keypoint sets that
respects their layout; and the two are combined by a fusion rule whose behaviour can be stated
exactly. Treating the stages as OT problems makes their design choices explicit and analysable: the
framework contains the assignments of NetVLAD and SALAD as special cases, its theory predicts when the
choice between them matters, and our experiments test those predictions with a frozen and with a
fine-tuned backbone. Our contributions are:
\begin{enumerate}
\item \textbf{An OT family of aggregators} (\cref{sec:uot}). A semi-relaxed unbalanced OT problem
defines token-to-cluster assignments indexed by one scalar $\lambda\in[0,1]$ that recovers the NetVLAD
soft-assignment ($\lambda=0$) and the balanced Sinkhorn assignment of SALAD~\cite{izquierdo2024optimal}
($\lambda=1$) \emph{exactly}; Galatea can also learn $\lambda$. We prove linear convergence with rate
$\lambda$ for the truncated iterations, characterise the dust-bin as a logistic gate on cluster
affinity, and show that intra-normalised descriptors see the transport plan only through per-token
weights (\cref{prop:uot,prop:contraction,prop:dustbin,prop:colinv}).
\item \textbf{OT-based global--local fusion} (\cref{sec:fgw,sec:fusion}). A training-free partial
fused Gromov--Wasserstein (FGW) matcher scores both the appearance and the relative layout of keypoints
and is invariant to similarity transforms (\cref{prop:inv}). A saturating fusion rule combines it with
the global similarity and provably overturns the global ranking only where its margin is small
(\cref{prop:override}).
\item \textbf{LunarSim-VPR} (\cref{sec:bench})\footnote{Code and data: \url{https://github.com/RaunakDune/galatea-lunar-vpr}. Project website and interactive demo: \url{https://lunarsim.org}.}, an open, procedurally generated benchmark
with a physically motivated renderer (crater size--frequency statistics, lunar-Lambert photometry
with opposition surge, cast shadows from local and regional relief) whose protocol separates
viewpoint change from controlled changes in solar elevation and azimuth.
\item \textbf{A controlled study of what matters} (\cref{sec:exp,sec:discussion}). Under identical
training, with three seeds and paired-bootstrap confidence intervals, we evaluate terrestrial models
zero-shot, the whole assignment family with a frozen and a fine-tuned backbone, four local scorers with
and without fusion, and the effect of training illumination. We distil the results into design
guidelines for lunar VPR.
\end{enumerate}

The findings are not all in favour of the most flexible model. With a frozen backbone every value of
$\lambda$ performs alike, exactly as \cref{prop:colinv} predicts. With a fine-tuned backbone the
family spreads only slightly, a learned $\lambda$ is no better than fixing it at its initial value, and
residual NetVLAD aggregation, a different design axis, pulls ahead.
Fusion, in contrast, helps every stage-1 model, and how local and global evidence are combined matters
more than which local matcher is used.

% ======================================================================
\section{Related Work}\label{sec:related}
% ======================================================================
\paragraph*{Global descriptors for VPR}
VLAD~\cite{jegou2010vlad} and its trainable counterpart NetVLAD~\cite{arandjelovic2016netvlad}
softly assign local features to clusters and sum their residuals. GeM pooling~\cite{radenovic2019gem}
and classification-based training at scale (CosPlace~\cite{berton2022cosplace},
EigenPlaces~\cite{berton2023eigenplaces}) give strong compact descriptors. MixVPR~\cite{alibey2023mixvpr}
mixes feature maps with MLPs. AnyLoc~\cite{keetha2024anyloc} showed that unsupervised aggregation of
DINOv2~\cite{oquab2024dinov2} features generalises across domains, including aerial and underwater
imagery. SALAD~\cite{izquierdo2024optimal} recast NetVLAD's assignment as balanced entropic OT
with a dust-bin on top of a partially fine-tuned DINOv2 backbone. Our aggregator generalises SALAD
to unbalanced OT and contains NetVLAD's assignment as a limiting case.

\paragraph*{Re-ranking and matching}
Two-stage pipelines re-rank global candidates with local evidence, for example multi-scale patch
descriptors~\cite{hausler2021patch} or keypoint matching with geometric verification. Learned
matchers such as SuperGlue~\cite{sarlin2020superglue} and LightGlue~\cite{lindenberger2023lightglue}
solve (or approximate) an OT problem with dust-bins over attention-refined
SuperPoint~\cite{detone2018superpoint} features. They are trained on terrestrial imagery and assume
that the descriptors carry enough appearance information. Our re-ranker is training-free. It adds a
Gromov--Wasserstein structural term~\cite{memoli2011gromov,peyre2016gromov}, in the fused form of
\cite{vayer2019fgw} and with partial mass~\cite{chapel2020partial}, so that matches must preserve
relative keypoint geometry.

\paragraph*{Optimal transport}
Entropic regularisation and the Sinkhorn algorithm~\cite{cuturi2013sinkhorn,peyre2019computational}
made OT differentiable and fast. Unbalanced OT replaces hard marginal constraints with divergence
penalties and keeps Sinkhorn-like scaling algorithms~\cite{chizat2018scaling}. Gromov--Wasserstein
distances compare metric spaces through their internal distances~\cite{memoli2011gromov} and are
used for structured matching~\cite{vayer2019fgw}. We use OT in both roles: as a layer inside the
descriptor and as a training-free matcher between keypoint sets.

\paragraph*{Planetary localisation}
Rovers rely on visual odometry~\cite{maimone2007mer}, and landers on terrain-relative navigation
against orbital maps~\cite{johnson2008trn}. Synthetic optical images rendered from DTMs support lunar
navigation~\cite{ceresoli2025optical}. Our DESPINA~\cite{nelsonarcher2026despina} synthesises
photorealistic horizon views for a given location and heading with DEM-conditioned diffusion.
Lunar crater images have been retrieved with visual and depth features~\cite{lunarcomplexfeatures}.
Surface-level place recognition under controlled illumination change remains largely unstudied.
